%! Principles of Experimental Design — Unit 1, Lesson 1.6 — Homework
\documentclass[10pt]{article}
\usepackage{saar-article}
\usepackage{saar-boxes}
\newcommand{\TallMath}[1]{$\displaystyle #1\rule[-1.4em]{0pt}{3.2em}$}

\begin{document}

\pageheader{Unit 1, Lesson 1.6}{Homework}
% NAMESTRIP: no \namedateperiod here. The student writes their name once, on the
% cover this component is stapled behind. See references/conventions.md.

\vspace{0.15in}

\boxguard[24]
\begin{scenariobox}[Millbrook Public Library --- This Time They Assign]{cerulean}
\small
Last summer Millbrook \emph{watched}: of the card holders who had signed
themselves up for the Summer Reading Program, $80\%$ read at least $10$ books,
against $32\%$ of those who had not --- a $48$-point gap the director could not
explain, because the card holders had sorted themselves.

\vspace{2pt}
This summer the library has $1{,}200$ card holders and runs an experiment.
\textbf{$200$ card holders volunteer.} A random number generator assigns
\textbf{$100$ to the Summer Reading Program} and \textbf{$100$ to a waiting
list} that starts in the fall. Everyone keeps the same library card and the same
borrowing limits. In September, \textbf{$65$} of the program group and
\textbf{$35$} of the waiting-list group had read at least $10$ books.
\end{scenariobox}

\vspace{0.08in}

\boxguard
\begin{notesbox}{Practice}
\small
\begin{enumerate}[leftmargin=1.6em, itemsep=8pt, topsep=3pt]

  \item Nobody let a card holder pick a group this year.

        \textbf{(a)} Is this an experiment? \blank{2cm}

        \textbf{(b)} Who decided which group each volunteer was in?
        \blank{5cm}

  \item Compute the reading rate for each group.

        \textbf{(a)} What percent of the $100$ in the program read at least
        $10$ books?

        \begin{work}
          \dfrac{65}{100} &= 0.65 \\
                          &= 65\%
        \end{work}

        \textbf{(b)} What percent of the $100$ on the waiting list did?

        \begin{work}
          \dfrac{35}{100} &= 0.35 \\
                          &= 35\%
        \end{work}

        \textbf{(c)} The program group is higher by \blank{2cm} percentage
        points.

  \item Now use all $200$ volunteers.

        \textbf{(a)} How many read at least $10$ books, and what percent is
        that?

        \begin{work}
          65 + 35 &= 100 \text{ card holders} \\
          \dfrac{100}{200} &= 0.50 = 50\%
        \end{work}

        \textbf{(b)} About how many of all $1{,}200$ card holders does that
        predict?

        \begin{work}
          0.50 \times 1200 &= 600 \text{ card holders}
        \end{work}

  \item Name the parts of this design.

        Experimental units: \blank{4cm} \hfill Treatment: \blank{4cm}

        Control group: \blank{4cm} \hfill Response variable: \blank{4cm}

  \item In a few words each, say how this design meets the four principles.

        Comparison: \blank{3.6cm} \hfill Randomization: \blank{3.6cm}

        Replication: \blank{3.6cm} \hfill Control: \blank{3.6cm}

  \tcbbreak   % keep item 6's prompt with its table — mirrored in homework_key
  \item Each row is a different plan the director could have used instead. Name
        the one principle it breaks.

        \begin{center}
        \renewcommand{\arraystretch}{1.45}
        \begin{tabularx}{0.97\linewidth}{@{} X p{3.7cm} @{}}
        \toprule
        \textbf{The plan} & \textbf{Principle broken} \\
        \midrule
        Enroll all $200$ volunteers in the program and see how much they read.
          & \blank{3.3cm} \\
        Enroll the first $100$ volunteers who signed the sheet; the rest wait.
          & \blank{3.3cm} \\
        Enroll $1$ volunteer and put $1$ on the waiting list.
          & \blank{3.3cm} \\
        Enroll the $100$ volunteers who live near the library; the $100$ who
        live far away wait. & \blank{3.3cm} \\
        \bottomrule
        \end{tabularx}
        \end{center}

  \item Of the $200$ volunteers, $120$ are adults and $80$ are children. The
        director decides to block on that and randomize inside each block.

        \begin{center}
        \renewcommand{\arraystretch}{1.4}
        \begin{tabularx}{0.9\linewidth}{@{} X c c c @{}}
        \toprule
        \textbf{Block} & \textbf{In block} & \textbf{To program}
          & \textbf{To waiting list} \\
        \midrule
        Adults & $120$ & \blank{1.5cm} & \blank{1.5cm} \\
        Children & $80$ & \blank{1.5cm} & \blank{1.5cm} \\
        \midrule
        Total & $200$ & \blank{1.5cm} & \blank{1.5cm} \\
        \bottomrule
        \end{tabularx}
        \end{center}

        Why would the director block on age rather than leave it to chance?
        \writeline

  \item Last summer's study found a $48$-point gap; this summer's experiment
        finds only $30$. Explain why the smaller gap is the \emph{stronger}
        evidence that the program works.
        \writelines{2}

\end{enumerate}
\end{notesbox}

\vspace{0.08in}

\boxguard
\begin{extensionbox}
\small
Take a headline you have seen that claims one thing causes another --- yours
from last lesson's extension will do. Design the experiment that would actually
test it:

\begin{enumerate}[leftmargin=1.6em, itemsep=3pt, topsep=3pt]
  \item Who or what are the experimental units, and how many would you need?
  \item What is the treatment, and what does the control group get instead?
  \item How exactly would chance assign the units to the groups?
  \item Name one variable you would hold the same for both groups, and say what
        would go wrong if you did not.
  \item Could anyone be blinded? Say who, or say why nobody could be.
\end{enumerate}
\end{extensionbox}

\vspace{0.08in}

\begin{remindbox}
\small
\textbf{Next class (Lesson 1.7):} both kinds of study side by side. When is an
experiment worth the cost, and when is it impossible or wrong to run one? You
will compare what an experiment and an observational study each let you conclude
and choose the data collection method that fits a given context.
\end{remindbox}

\end{document}
